% Generated by scripts/assemble.py; corrections belong in source/corrections.json.
\examhead{32}{首都师范大学 2025 年数学分析试卷}{35}
% p036-b17; source PDF page 36
\par\Needspace{5\baselineskip}
\blockanchor{p036-b17}
\majorquestion{一}(24 分) 计算题.\par
% p036-b18; source PDF page 36
\blockanchor{p036-b18}
\question{1}(8 分) 求极限 $\lim_{x \to 0} \left( \frac{\sin x}{x} \right)^{\frac{1}{1 - \cos x}}$.\par
% p036-b19; source PDF page 36
\blockanchor{p036-b19}
\question{2}(8 分) 设函数 $f(x)$ 在点 $a$ 可导且 $f(a) \neq 0$, 求 $\lim_{n \to \infty} \left[ \frac{f\left(a + \frac{1}{n}\right)}{f(a)} \right]^n$.\par
% p036-b20; source PDF page 36
\blockanchor{p036-b20}
\question{3}(8 分) 求级数 $\sum_{n=1}^{\infty} \frac{n}{3^n}$ 的和.\par
% p036-b21; source PDF page 36
\par\Needspace{5\baselineskip}
\par\Needspace{6\baselineskip}
\blockanchor{p036-b21}
\majorquestion{二}(15 分) 考虑方程\par
\nopagebreak[4]
% p036-b22; source PDF page 36
\blockanchor{p036-b22}
\displayformula{e ^ {x + 2 y + 3 z} + x y z = 1.}
% p036-b23; source PDF page 36
\blockanchor{p036-b23}
\question{1}(7 分) 设函数 $z = z(x, y)$ 由 (??) 确定, 求 $\mathrm{dz}$.\par
% p036-b24; source PDF page 36
\blockanchor{p036-b24}
\question{2}(8 分) 设曲面 $S$ 由 (??) 确定, 求 $S$ 上过点 $(0,0,0)$ 的切平面方程和法线方程.\par
% p036-b25; source PDF page 36
\par\Needspace{5\baselineskip}
\blockanchor{p036-b25}
\majorquestion{三}(16 分) 考虑函数 $f(x) = \begin{cases} \frac{x}{3} + x^2 \sin \frac{1}{x}, & x \neq 0; \\ 0, & x = 0. \end{cases}$\par
% p036-b26; source PDF page 36
\blockanchor{p036-b26}
\question{1}(8 分) 讨论函数 $f(x)$ 的连续性和可导性.\par
% p036-b27; source PDF page 36
\blockanchor{p036-b27}
\question{2}(8 分) 讨论是否存在 $x = 0$ 的一个邻域, 使得 $f(x)$ 在该邻域内单调并给出理由.\par
% p037-b00; source PDF page 37
\par\Needspace{5\baselineskip}
\blockanchor{p037-b00}
\majorquestion{四}(10 分) 证明不等式 $xy \leq x \ln x - x + e^y$, $x \geq 1$, $y \geq 0$.\par
% p037-b01; source PDF page 37
\par\Needspace{5\baselineskip}
\blockanchor{p037-b01}
\majorquestion{五}(12 分) 设 $L$ 为球面 $x^{2} + y^{2} + z^{2} = a^{2}$ ($a > 0$) 和平面 $x + y + z = 0$ 的交线, 一质点在点 $M(x, y, z)$ 处受到力 $\pmb{F} = y\pmb{i} + z\pmb{j} + x\pmb{k}$ 的作用, 其中 $\pmb{i} = (1, 0, 0), \pmb{j} = (0, 1, 0), \pmb{k} = (0, 0, 1)$, 此质点沿 $L$ 运动一周, 方向从 $z$ 轴正项看为逆时针方向, 求 $\pmb{F}$ 所作的功 $W$.\par
% p037-b02; source PDF page 37
\par\Needspace{5\baselineskip}
\blockanchor{p037-b02}
\majorquestion{六}(10 分) 设  $f(x)$  是区间  $[a, b]$  上的连续函数, 且满足条件  $\int_{a}^{b} x^{k} f(x) \, dx = 0, k = 0, 1, \cdots, n$ . 证明: 函数  $f(x)$  在  $(a, b)$  内至少有  $n + 1$  个不同的零点.\par
% p037-b03; source PDF page 37
\par\Needspace{5\baselineskip}
\blockanchor{p037-b03}
\majorquestion{七}(20 分) 设 $f(x) = \sum_{n=1}^{\infty} \frac{\cos nx}{n^2 + 1}$.\par
% p037-b04; source PDF page 37
\blockanchor{p037-b04}
\question{1}(10 分) 证明: $f(x)$ 在区间 $(0,2\pi)$ 内连续, 且有连续的导函数.\par
% p037-b05; source PDF page 37
\blockanchor{p037-b05}
\question{2}(10 分) 记 $F(x) = \int_{0}^{x} f(t) \, \mathrm{d}t$, 证明: $-\frac{1}{30} < F\left(\frac{\pi}{2}\right) < \frac{1}{2}$.\par
% p037-b06; source PDF page 37
\par\Needspace{5\baselineskip}
\par\Needspace{6\baselineskip}
\blockanchor{p037-b06}
\majorquestion{八}(18 分) 讨论下列反常积分的敛散性(包括绝对收敛, 条件收敛和发散)\par
\nopagebreak[4]
% p037-b07; source PDF page 37
\blockanchor{p037-b07}
\displayformula{\int_ {0} ^ {+ \infty} \frac {x \sin x}{1 + x ^ {p}} \mathrm{d} x, p \geq 0.}
% p037-b08; source PDF page 37
\par\Needspace{5\baselineskip}
\blockanchor{p037-b08}
\majorquestion{九}(15 分) 设 $f(x)$ 以 $2\pi$ 为周期且具有一阶连续的导函数. 证明: $f(x)$ 的傅里叶级数在 $(-\infty, +\infty)$ 上一致收敛于 $f(x)$.\par
% p037-b09; source PDF page 37
\par\Needspace{5\baselineskip}
\blockanchor{p037-b09}
\majorquestion{十}(10 分) (本题可能有误) 若函数 $f(x)$ 在 (0,1) 内可微, 且 $\lim_{x\to 0^{+}}f(x) = \infty$, 试证: $\lim_{x\to 0^{+}}f'(x) = \infty$.\par
